\begin{theorem}[Literal weighted contractions under vertex gauges]
    \label{thm:peps_regular_weighted_gauge_transport}
    \lean{TNLean.PEPS.graphInsertedBondNetwork_diagonal_vertexGauge,
        TNLean.PEPS.regularWeightedOpenRegionWeight_vertexGauge,
        TNLean.PEPS.regularWeightedOpenRegionWeight_chargePair_swept}
    \leanok
    \uses{thm:peps_regular_weighted_bond_contraction,
        def:peps_regular_weighted_open_contraction,
        def:peps_regular_swept_patch}
    Let $G$ be finite, and let the original site coefficients be invariant
    under simultaneous left translation of their incident labels. Orient each
    edge from its smaller endpoint $t(e)$ to its larger endpoint $h(e)$.
    For a vertex gauge $\ell$, put
    \[
        u'_e=\ell_{h(e)}^{-1}u_e\ell_{t(e)},\qquad
        (L\eta)_e=\ell_{t(e)}\eta_e,\qquad
        \theta'_e=\ell_{t(e)}^{-1}\theta_e.
    \]
    The actual weighted open coefficient satisfies
    \[
        \Psi_R(u,F;\theta)=\Psi_R(u',F\circ L;\theta').
    \]
    The weight $F$ may correlate several edges. For independent diagonal
    edge weights, the same equality holds for the actual closed contraction
    with bond matrices $U_{u_e}\operatorname{diag}(f_e)$; the new weight is
    $f'_e(t)=f_e(\ell_{t(e)}t)$.

    In particular, choose $\ell=k^{-1}$ on a swept set $S$ and $1$ outside.
    If the first selected edge has its tail in $S$ and the second does not,
    the correlated weight transforms as
    \[
        \chi(p\eta_1^{-1}\eta_0)
        \longmapsto\chi(p\eta_1^{-1}k^{-1}\eta_0).
    \]
    These are auxiliary actual contraction identities for
    \cite{Schuch2010PEPS}, lines 2505--2535 and 2560--2581.
    No expansion of character diagonals into group matrices is required.
\end{theorem}
\begin{proof}\leanok
    Change each summed tail label to $\ell_{t(e)}^{-1}\eta_e$.
    The new head label is $\ell_{h(e)}^{-1}u_e\eta_e$.
    Local translation invariance cancels the vertex factors, and the
    restriction of this bijection to crossing edges gives $\theta'$.
    The scalar weight is retained in the same finite sum.
\end{proof}

\begin{definition}[The charge edge in the native string-deformation figure]
    \label{def:peps_torus_charge_string_edge}
    \lean{TNLean.PEPS.torusSweptStringChargeEdge}
    \leanok
    \uses{def:peps_torus_swept_string_patch}
    At a translated reference vertex $v=(x,y)$, let $e_c$ be the native
    horizontal bond joining $(x+1,y+1)$ to $(x+2,y+1)$.
    This is the central charge diagonal in
    \cite{Schuch2010PEPS}, lines 2569--2581 and the virtual charge--flux
    deformation figure. Its orientation is the actual ordered edge orientation.
\end{definition}

\begin{theorem}[The literal native charge--flux deformation]
    \label{thm:peps_torus_charge_string_deformation}
    \lean{TNLean.PEPS.torusSweptStringChargeEdge_tail_mem,
        TNLean.PEPS.IsGIsometric.graphInsertedBondNetwork_torusChargeStringDeformation}
    \leanok
    \uses{thm:peps_regular_weighted_gauge_transport,
        def:peps_torus_charge_string_edge,
        thm:peps_torus_swept_string_transport}
    Let the torus periods satisfy $w\geq5$ and $H\geq4$, and let the native
    four-leg tensor be regular $G$-isometric. Use the literal initial flux
    string $u^{(0)}$ with label $k$ and no second flux string. Insert the
    diagonal $\operatorname{diag}(\chi(pt))$ on $e_c$ and unit weights
    on the other bonds. Its actual original-spin contraction equals the
    contraction with the swept string $u^{(j)}$, for either displayed
    sweep $j=1,2$, and with diagonal
    \[
        \operatorname{diag}\bigl(\chi(pk^{-1}t)\bigr)
    \]
    on the same charge edge. This holds at every translated position,
    including both periodic seams, and for every scalar function $\chi$.
    When $\chi$ is a character its type is unchanged.

    This is the actual gauge-presentation identity of the virtual figure
    in \cite{Schuch2010PEPS}, lines 2560--2581. It does not by itself
    identify a physical crossing route or a reunion experiment with a braid.
\end{theorem}
\begin{proof}\leanok
    Both endpoints of $e_c$ belong to each displayed swept set.
    Thus its tail gauge is $k^{-1}$ in either native orientation.
    Native regular $G$-isometry supplies local translation invariance.
    Apply the weighted gauge identity with all other diagonal weights equal
    to one and the actual swept bond assignment.
\end{proof}
